% --- Archon physics formalization source begin ---
% archon:physics
% archon:covers PhyXMiniProblems/problem_phyx_mini_0289.lean
% archon:source-report reports/phyx_mini/problem_phyx_mini_0289.source.json

\chapter{Physics problem phyx\_mini\_0289}
\label{ch:PhyXMiniProblems_problem_phyx_mini_0289}

\paragraph{Problem source.}
\#\# Physical scenario

A sinusoidal wave moving along a string is shown twice in figure, as crest \$A\$ travels in the positive direction of an \$x\$ axis by distance \$d = 6.0\$ cm in \$4.0\$ ms. The tick marks along the axis are separated by \$10\$ cm; height \$H = 6.00\$ mm.

\#\# Question

The equation for the wave is in the form \$y(x, t) = y\_m \textbackslash{}sin(kx \textbackslash{}pm \textbackslash{}omega t)\$, so what is \$\textbackslash{}omega\$?

\#\# Answer choices

- A: A: \$1.8\textbackslash{}times 10\textasciicircum{}\{2\}\$ rad/s

- B: B: \$2.0\textbackslash{}times 10\textasciicircum{}\{2\}\$ rad/s

- C: C: \$2.8\textbackslash{}times 10\textasciicircum{}\{2\}\$ rad/s

- D: D: \$2.4\textbackslash{}times 10\textasciicircum{}\{2\}\$ rad/s

\#\# Figure caption (auxiliary; use the image as primary evidence)

The image depicts two sinusoidal waveforms on an x-y graph. 

1. **Waveforms**:
   - The solid line represents one wave and the dashed line represents another.
   - Both waves have a similar shape and amplitude, showing periodic oscillation.

2. **Axes and Labels**:
   - The x-axis is labeled "x."
   - The y-axis is labeled "y."
   - "H" is labeled vertically along the y-axis, indicating the height or amplitude of the waves.
   - "A" is labeled at the peak of the solid wave, likely indicating the amplitude.
   - "d" is labeled with a horizontal arrow between the waves, indicating the phase difference or distance.

3. **Lines and Arrows**:
   - There is a horizontal line along the x-axis acting as the baseline for the waves.
   - The arrows and double-ended arrows indicate measurements (height and distance).
  
These elements combined suggest an illustration of wave interference or phase shift.

\#\# Dataset metadata

Wave Properties; Physical Model Grounding Reasoning, Numerical Reasoning

\paragraph{Recorded answer/context.}
D: D: \$2.4\textbackslash{}times 10\textasciicircum{}\{2\}\$ rad/s

\paragraph{Figure/image path.}
/root/proposal\_for\_physic/science-mango/phyx\_mini\_run/phyx\_data/test\_image/289.png

\paragraph{Formalization target.}
create a compiling Lean file with sorry bodies at `PhyXMiniProblems/problem\_phyx\_mini\_0289.lean`. The Lean declarations must preserve the physical quantities, dimensions or dimensional roles, figure labels, governing-law hypotheses, and final relation expressed by this problem.
Use Mathlib/Physlib names found through LeanExplore where available. If a physics API is missing, introduce faithful local abstractions rather than scalar placeholder aliases.

\begin{theorem}[Physics formalization target]
\label{thm:physics:phyx_mini_0289:target}
\lean{PhyXMiniProblems.ProblemPhyXMini0289.problem_phyx_mini_0289}
\uses{def:physics:phyx-mini-0289:phyxminiproblems-problemphyxmini0289-angularfrequencyinradianspersecond, def:physics:phyx-mini-0289:phyxminiproblems-problemphyxmini0289-travelingstringwaveexperiment, def:physics:phyx-mini-0289:phyxminiproblems-problemphyxmini0289-matchesproblemandprimaryfigure, def:physics:phyx-mini-0289:phyxminiproblems-problemphyxmini0289-hasphysicalwaveparameters, def:physics:phyx-mini-0289:phyxminiproblems-problemphyxmini0289-satisfiestravelingsinusoidalwavelaws, def:physics:phyx-mini-0289:phyxminiproblems-problemphyxmini0289-matchesnearesttendisplay}
The assigned autoformalize agent should translate this physics problem into Lean declarations in the covered file, with theorem and lemma proof bodies written as `by sorry`.
\end{theorem}
\begin{proof}
This is an autoformalization task, not a proof task. Produce statements that can later be proved by the physics prover without weakening the physical model.
\end{proof}

% --- Archon declaration topology begin ---
\section{Declaration topology}
\begin{definition}[Length Quantity]
  \label{def:physics:phyx-mini-0289:phyxminiproblems-problemphyxmini0289-lengthquantity}
  \lean{PhyXMiniProblems.ProblemPhyXMini0289.LengthQuantity}
  A nonnegative physical length
\end{definition}
\begin{proof}
  This is the declared model definition of Length Quantity.
\end{proof}
\begin{definition}[Time Quantity]
  \label{def:physics:phyx-mini-0289:phyxminiproblems-problemphyxmini0289-timequantity}
  \lean{PhyXMiniProblems.ProblemPhyXMini0289.TimeQuantity}
  A nonnegative physical duration
\end{definition}
\begin{proof}
  This is the declared model definition of Time Quantity.
\end{proof}
\begin{definition}[Speed Quantity]
  \label{def:physics:phyx-mini-0289:phyxminiproblems-problemphyxmini0289-speedquantity}
  \lean{PhyXMiniProblems.ProblemPhyXMini0289.SpeedQuantity}
  A nonnegative propagation speed along the string
\end{definition}
\begin{proof}
  This is the declared model definition of Speed Quantity.
\end{proof}
\begin{definition}[Wave Number Quantity]
  \label{def:physics:phyx-mini-0289:phyxminiproblems-problemphyxmini0289-wavenumberquantity}
  \lean{PhyXMiniProblems.ProblemPhyXMini0289.WaveNumberQuantity}
  A nonnegative wave number; radians are dimensionless
\end{definition}
\begin{proof}
  This is the declared model definition of Wave Number Quantity.
\end{proof}
\begin{definition}[Angular Frequency Quantity]
  \label{def:physics:phyx-mini-0289:phyxminiproblems-problemphyxmini0289-angularfrequencyquantity}
  \lean{PhyXMiniProblems.ProblemPhyXMini0289.AngularFrequencyQuantity}
  A nonnegative angular frequency; radians are dimensionless
\end{definition}
\begin{proof}
  This is the declared model definition of Angular Frequency Quantity.
\end{proof}
\begin{definition}[length Readout]
  \label{def:physics:phyx-mini-0289:phyxminiproblems-problemphyxmini0289-lengthreadout}
  \lean{PhyXMiniProblems.ProblemPhyXMini0289.lengthReadout}
  \uses{def:physics:phyx-mini-0289:phyxminiproblems-problemphyxmini0289-lengthquantity}
  Numerical readout of a physical length in the selected length unit
\end{definition}
\begin{proof}
  This is the declared model definition of length Readout.
\end{proof}
\begin{definition}[time Readout]
  \label{def:physics:phyx-mini-0289:phyxminiproblems-problemphyxmini0289-timereadout}
  \lean{PhyXMiniProblems.ProblemPhyXMini0289.timeReadout}
  \uses{def:physics:phyx-mini-0289:phyxminiproblems-problemphyxmini0289-timequantity}
  Numerical readout of a duration in the selected time unit
\end{definition}
\begin{proof}
  This is the declared model definition of time Readout.
\end{proof}
\begin{definition}[speed Readout]
  \label{def:physics:phyx-mini-0289:phyxminiproblems-problemphyxmini0289-speedreadout}
  \lean{PhyXMiniProblems.ProblemPhyXMini0289.speedReadout}
  \uses{def:physics:phyx-mini-0289:phyxminiproblems-problemphyxmini0289-speedquantity}
  Numerical readout of a speed in the selected length and time units
\end{definition}
\begin{proof}
  This is the declared model definition of speed Readout.
\end{proof}
\begin{definition}[wave Number In Radians Per Meter]
  \label{def:physics:phyx-mini-0289:phyxminiproblems-problemphyxmini0289-wavenumberinradianspermeter}
  \lean{PhyXMiniProblems.ProblemPhyXMini0289.waveNumberInRadiansPerMeter}
  \uses{def:physics:phyx-mini-0289:phyxminiproblems-problemphyxmini0289-wavenumberquantity}
  Radian-per-meter readout of a wave number
\end{definition}
\begin{proof}
  This is the declared model definition of wave Number In Radians Per Meter.
\end{proof}
\begin{definition}[angular Frequency In Radians Per Second]
  \label{def:physics:phyx-mini-0289:phyxminiproblems-problemphyxmini0289-angularfrequencyinradianspersecond}
  \lean{PhyXMiniProblems.ProblemPhyXMini0289.angularFrequencyInRadiansPerSecond}
  \uses{def:physics:phyx-mini-0289:phyxminiproblems-problemphyxmini0289-angularfrequencyquantity}
  Radian-per-second readout of an angular frequency
\end{definition}
\begin{proof}
  This is the declared model definition of angular Frequency In Radians Per Second.
\end{proof}
\begin{definition}[Snapshot]
  \label{def:physics:phyx-mini-0289:phyxminiproblems-problemphyxmini0289-snapshot}
  \lean{PhyXMiniProblems.ProblemPhyXMini0289.Snapshot}
  The two depictions of the same wave in the primary figure
\end{definition}
\begin{proof}
  This is the declared model definition of Snapshot.
\end{proof}
\begin{definition}[Curve Style]
  \label{def:physics:phyx-mini-0289:phyxminiproblems-problemphyxmini0289-curvestyle}
  \lean{PhyXMiniProblems.ProblemPhyXMini0289.CurveStyle}
  Line styles distinguishing the two snapshots
\end{definition}
\begin{proof}
  This is the declared model definition of Curve Style.
\end{proof}
\begin{definition}[Axis]
  \label{def:physics:phyx-mini-0289:phyxminiproblems-problemphyxmini0289-axis}
  \lean{PhyXMiniProblems.ProblemPhyXMini0289.Axis}
  The coordinate axes explicitly labelled in the figure
\end{definition}
\begin{proof}
  This is the declared model definition of Axis.
\end{proof}
\begin{definition}[Axis Label]
  \label{def:physics:phyx-mini-0289:phyxminiproblems-problemphyxmini0289-axislabel}
  \lean{PhyXMiniProblems.ProblemPhyXMini0289.AxisLabel}
  Printed axis labels
\end{definition}
\begin{proof}
  This is the declared model definition of Axis Label.
\end{proof}
\begin{definition}[Figure Role]
  \label{def:physics:phyx-mini-0289:phyxminiproblems-problemphyxmini0289-figurerole}
  \lean{PhyXMiniProblems.ProblemPhyXMini0289.FigureRole}
  Roles of the three letter labels visible in the figure
\end{definition}
\begin{proof}
  This is the declared model definition of Figure Role.
\end{proof}
\begin{definition}[Figure Label]
  \label{def:physics:phyx-mini-0289:phyxminiproblems-problemphyxmini0289-figurelabel}
  \lean{PhyXMiniProblems.ProblemPhyXMini0289.FigureLabel}
  The labels printed beside the crest and the two measured spans
\end{definition}
\begin{proof}
  This is the declared model definition of Figure Label.
\end{proof}
\begin{definition}[Propagation Direction]
  \label{def:physics:phyx-mini-0289:phyxminiproblems-problemphyxmini0289-propagationdirection}
  \lean{PhyXMiniProblems.ProblemPhyXMini0289.PropagationDirection}
  Direction of propagation along the labelled horizontal axis
\end{definition}
\begin{proof}
  This is the declared model definition of Propagation Direction.
\end{proof}
\begin{definition}[Traveling String Wave Experiment]
  \label{def:physics:phyx-mini-0289:phyxminiproblems-problemphyxmini0289-travelingstringwaveexperiment}
  \lean{PhyXMiniProblems.ProblemPhyXMini0289.TravelingStringWaveExperiment}
  \uses{def:physics:phyx-mini-0289:phyxminiproblems-problemphyxmini0289-lengthquantity, def:physics:phyx-mini-0289:phyxminiproblems-problemphyxmini0289-timequantity, def:physics:phyx-mini-0289:phyxminiproblems-problemphyxmini0289-speedquantity, def:physics:phyx-mini-0289:phyxminiproblems-problemphyxmini0289-wavenumberquantity, def:physics:phyx-mini-0289:phyxminiproblems-problemphyxmini0289-angularfrequencyquantity, def:physics:phyx-mini-0289:phyxminiproblems-problemphyxmini0289-snapshot, def:physics:phyx-mini-0289:phyxminiproblems-problemphyxmini0289-curvestyle, def:physics:phyx-mini-0289:phyxminiproblems-problemphyxmini0289-axis, def:physics:phyx-mini-0289:phyxminiproblems-problemphyxmini0289-axislabel, def:physics:phyx-mini-0289:phyxminiproblems-problemphyxmini0289-figurerole, def:physics:phyx-mini-0289:phyxminiproblems-problemphyxmini0289-figurelabel, def:physics:phyx-mini-0289:phyxminiproblems-problemphyxmini0289-propagationdirection}
  Independent physical quantities and scalar field for the string wave. displacementInMeters x t is the signed transverse displacement in meters at an SI position x and SI time t. It is a scalar component readout, not a replacement for any of the dimensionful physical quantities
\end{definition}
\begin{proof}
  This is the declared model definition of Traveling String Wave Experiment.
\end{proof}
\begin{definition}[Matches Problem And Primary Figure]
  \label{def:physics:phyx-mini-0289:phyxminiproblems-problemphyxmini0289-matchesproblemandprimaryfigure}
  \lean{PhyXMiniProblems.ProblemPhyXMini0289.MatchesProblemAndPrimaryFigure}
  \uses{def:physics:phyx-mini-0289:phyxminiproblems-problemphyxmini0289-lengthreadout, def:physics:phyx-mini-0289:phyxminiproblems-problemphyxmini0289-timereadout, def:physics:phyx-mini-0289:phyxminiproblems-problemphyxmini0289-travelingstringwaveexperiment}
  Problem data and geometric readouts from the primary figure. The first four fields transcribe the stated numerical data. The remaining fields record that the solid curve is earlier, the dashed curve is later, crest A moves in positive x, the labelled crest is at the fifth tick in the earlier snapshot, a wavelength spans four tick intervals, and H spans twice the amplitude. None constrains the requested angular frequency
\end{definition}
\begin{proof}
  This is the declared model definition of Matches Problem And Primary Figure.
\end{proof}
\begin{definition}[Has Physical Wave Parameters]
  \label{def:physics:phyx-mini-0289:phyxminiproblems-problemphyxmini0289-hasphysicalwaveparameters}
  \lean{PhyXMiniProblems.ProblemPhyXMini0289.HasPhysicalWaveParameters}
  \uses{def:physics:phyx-mini-0289:phyxminiproblems-problemphyxmini0289-lengthreadout, def:physics:phyx-mini-0289:phyxminiproblems-problemphyxmini0289-timereadout, def:physics:phyx-mini-0289:phyxminiproblems-problemphyxmini0289-speedreadout, def:physics:phyx-mini-0289:phyxminiproblems-problemphyxmini0289-wavenumberinradianspermeter, def:physics:phyx-mini-0289:phyxminiproblems-problemphyxmini0289-angularfrequencyinradianspersecond, def:physics:phyx-mini-0289:phyxminiproblems-problemphyxmini0289-travelingstringwaveexperiment}
  Positivity assumptions selecting a nondegenerate traveling wave
\end{definition}
\begin{proof}
  This is the declared model definition of Has Physical Wave Parameters.
\end{proof}
\begin{definition}[Satisfies Traveling Sinusoidal Wave Laws]
  \label{def:physics:phyx-mini-0289:phyxminiproblems-problemphyxmini0289-satisfiestravelingsinusoidalwavelaws}
  \lean{PhyXMiniProblems.ProblemPhyXMini0289.SatisfiesTravelingSinusoidalWaveLaws}
  \uses{def:physics:phyx-mini-0289:phyxminiproblems-problemphyxmini0289-lengthreadout, def:physics:phyx-mini-0289:phyxminiproblems-problemphyxmini0289-timereadout, def:physics:phyx-mini-0289:phyxminiproblems-problemphyxmini0289-speedreadout, def:physics:phyx-mini-0289:phyxminiproblems-problemphyxmini0289-wavenumberinradianspermeter, def:physics:phyx-mini-0289:phyxminiproblems-problemphyxmini0289-angularfrequencyinradianspersecond, def:physics:phyx-mini-0289:phyxminiproblems-problemphyxmini0289-travelingstringwaveexperiment}
  The standard kinematic and sinusoidal-wave laws used by the problem: the translated crest obeys d = v Δt in every coherent pair of units; wave number and wavelength obey k λ = 2π; angular frequency obeys ω = k v; a wave traveling in positive x has phase kx - ωt. These are general physical laws and contain neither 75π nor any displayed answer value
\end{definition}
\begin{proof}
  This is the declared model definition of Satisfies Traveling Sinusoidal Wave Laws.
\end{proof}
\begin{definition}[Answer Choice]
  \label{def:physics:phyx-mini-0289:phyxminiproblems-problemphyxmini0289-answerchoice}
  \lean{PhyXMiniProblems.ProblemPhyXMini0289.AnswerChoice}
  Labels printed beside the four candidate angular frequencies
\end{definition}
\begin{proof}
  This is the declared model definition of Answer Choice.
\end{proof}
\begin{definition}[radians Per Second]
  \label{def:physics:phyx-mini-0289:phyxminiproblems-problemphyxmini0289-answerchoice-radianspersecond}
  \lean{PhyXMiniProblems.ProblemPhyXMini0289.AnswerChoice.radiansPerSecond}
  \uses{def:physics:phyx-mini-0289:phyxminiproblems-problemphyxmini0289-answerchoice}
  Radian-per-second value printed beside each answer label
\end{definition}
\begin{proof}
  This is the declared model definition of radians Per Second.
\end{proof}
\begin{definition}[recorded Answer Choice]
  \label{def:physics:phyx-mini-0289:phyxminiproblems-problemphyxmini0289-recordedanswerchoice}
  \lean{PhyXMiniProblems.ProblemPhyXMini0289.recordedAnswerChoice}
  \uses{def:physics:phyx-mini-0289:phyxminiproblems-problemphyxmini0289-answerchoice}
  The answer label recorded in the supplied dataset
\end{definition}
\begin{proof}
  This is the declared model definition of recorded Answer Choice.
\end{proof}
\begin{definition}[Matches Nearest Ten Display]
  \label{def:physics:phyx-mini-0289:phyxminiproblems-problemphyxmini0289-matchesnearesttendisplay}
  \lean{PhyXMiniProblems.ProblemPhyXMini0289.MatchesNearestTenDisplay}
  \uses{def:physics:phyx-mini-0289:phyxminiproblems-problemphyxmini0289-answerchoice}
  Agreement with a value displayed to the nearest ten radians per second. The strict half-width also excludes a midpoint tie
\end{definition}
\begin{proof}
  This is the declared model definition of Matches Nearest Ten Display.
\end{proof}
% --- Archon declaration topology end ---
% --- Archon physics formalization source end ---
